\documentclass{article}
\usepackage{neurips_2026}
\usepackage[utf8]{inputenc}
\usepackage[T1]{fontenc}
\usepackage{hyperref}
\usepackage{url}
\usepackage{booktabs}
\usepackage{amsfonts}
\usepackage{amsmath}
\usepackage{amssymb}
\usepackage{nicefrac}
\usepackage{microtype}
\usepackage{xcolor}
\usepackage{comment}
\usepackage{graphicx}

\title{A Graph-Based Inspection and Intervention Tool for \\
Validating Mechanistic Learning in PINNs}

\author{%
  Anonymous Author(s) \\
  Affiliation \\
  Address \\
  \texttt{email@domain.com}
}

\begin{document}
\maketitle

\begin{abstract}
We ask whether physically meaningful correspondences discovered inside a
trained scientific model remain meaningful outside the conditions under
which they were discovered. We introduce GIIT (Graph-based Inspection and
Intervention Tool), which represents governing physics as a causal graph,
maps graph nodes to internal network components via sensitivity- and
trend-based discovery, and tests the resulting mapping under targeted
intervention. Evaluating on temporal extrapolation out-of-distribution (OOD)
regimes, we find that extrapolating neural states suffer representation collapse and lose directional parameter decoupling. 
Specifically, while the model achieves $1.13\%$ $L_2$ error in-distribution, temporal extrapolation causes 
predictive error to explode ($>2000\%$), physical residuals to jump three orders of magnitude, and parameter sensitivities to invert sign, 
revealing mechanistic failure modes invisible to standard in-distribution metrics.

\vspace{0.5em}
\noindent\textbf{Keywords:} Physics-Informed Neural Networks, Mechanistic Interpretability, Scientific Machine Learning, Representation Analysis, Intervention-Based Interpretability, Out-of-Distribution Generalization, Physical Reasoning
\end{abstract}

\section{Introduction}
\label{sec:intro}
Physics-Informed Neural Networks (PINNs) \citep{raissi2019pinn} embed
differential-equation residuals directly into the training loss, achieving
strong performance in low-data, high-simulation-cost regimes. However, a
near-zero physics residual and low output error are a form of
\emph{numerical} constraint satisfaction, not a guarantee of
\emph{mechanistic} correctness: a trained PINN's hidden layers may encode a
representation with little correspondence to the causal structure of the
system it was fit to \citep{wang2024understanding,krishnapriyan2021characterizing}.

General frameworks
for understanding model behavior \citep{murdoch2019pdr} and XAI approaches
tailored to scientific settings \citep{roscher2020xai} both emphasize that
explanations should align with domain knowledge, but the dominant tools
remain correlational: saliency and SHAP-style attribution, and
counterfactual explanation methods \citep{chou2022counterfactual}, operate
on model \emph{inputs} rather than internal components. Causal approaches
such as Causal Dependence Plots \citep{loftus2023cdp} and DAG-based SHAP
\citep{terminassian2023dag} use causal structure to sharpen input-level
explanations, but likewise stop at the input-output boundary. Work that
does look inside the model tends to go the other direction (discovering
causal structure \emph{from} data or representations
\citep{renero2025rex}), whereas we start from known physical relationships
and test whether a trained model is consistent with them. Operator-learning
methods such as DeepONet \citep{lu2021deeponet} extend PINN-style training
to function-space mappings but do not test internal correspondence either.
Our framing draws on the role of causal structure in scientific explanation
\citep{buijsman2023causal} and on causal abstraction as a theoretical basis
for mechanistic interpretability \citep{geiger2025mechinterp}, applying
these ideas to PINNs via an explicit mapping between physical variables and
internal tensors, followed by targeted intervention. Recent work has used PINNs to quantify the sensitivity of physical solutions to governing parameters through automatic differentiation. Existing PINN sensitivity analysis primarily studies physical-output sensitivity to physical parameters \citep{hanna2024sensitivity}; GIIT studies sensitivity and intervention at the level of internal neural representations.

This motivates our central question: \emph{do physically meaningful
correspondences discovered inside a model remain meaningful when the model
is evaluated outside the conditions under which the correspondence was
found?} Specifically, we evaluate whether internal representations preserve
their mechanistic consistency when transitioning from in-distribution interpolation
($t \in [0, 1]$) to out-of-distribution temporal extrapolation ($t \in [1, 2]$).

\textbf{Contributions.} We make three key contributions:
\begin{itemize}
    \item \textbf{The GIIT Pipeline:} We introduce a tool that automatically links physical concepts (like mass or damping) to specific internal layers of a neural network using sensitivity analysis.
    \item \textbf{Intervention Testing:} We provide a way to test these links. By masking or adding noise to specific layers, we prove whether the network actually relies on them for physical reasoning.
    \item \textbf{Out-of-Distribution (OOD) Analysis:} We test whether these internal physical representations hold up under distribution shift. By freezing the mappings found during training and evaluating them on new data, we show when and how the model's physical understanding collapses.
\end{itemize}

\subsection{The GIIT Framework}
Figure~\ref{fig:overview} summarizes the pipeline: a physical specification
is compiled into a physics graph; a trained neural model is parsed for
candidate tensors; a variable/parameter mapping is discovered between graph
nodes and tensors; targeted interventions are applied to mapped tensors; and
the resulting change in physical observables is compared against the
analytically expected response.

\begin{figure}[h]
\centering
\includegraphics[width=0.85\linewidth]{figures/arch.jpg}
\caption{GIIT pipeline: Physical specification
$\rightarrow$ Physics graph $\rightarrow$ Neural model $\rightarrow$
Variable/parameter mapping $\rightarrow$ Intervention $\rightarrow$ Physical
observables $\rightarrow$ ID$\to$OOD transfer test.}
\label{fig:overview}
\end{figure}

\section{Method}
\label{sec:method}

\subsection{Physics Graph and Executable Ruleset}
We represent the governing physical mechanism as a directed graph
$G=(V,E)$, where $V$ is a set of physical variables (parameters, e.g. mass,
damping, viscosity; intermediate quantities, e.g. gradients; and outputs)
and $E$ encodes known directional dependencies (e.g. mass
$\rightarrow$ acceleration). Some nodes are direct model outputs; others are
derived via a small \emph{executable ruleset} on top of automatic
differentiation, e.g. $x \rightarrow x_t, x_{tt}$, or computed observables
such as amplitude $x \rightarrow A = \max(x) - \min(x)$. The graph can be
supplied directly as a JSON schema by a domain expert, or optionally
LLM-assisted from source code. The graph specifies \emph{what} physical
concepts exist; the ruleset specifies \emph{how} derived observables are
computed from model output.

\subsection{Neural-Physical Mapping}
GIIT uses two complementary discovery strategies to map a causal node $n$
to a concrete internal tensor $\ell^\ast$.

\textbf{Variables (sensitivity-based gradient mapping).} For a node
corresponding to a quantity computable from model output via a known
formula $F_n$, we compute $q = \text{mean}(F_n(y))$ for $y=M(x)$ and, for
every candidate layer $\ell$,
\begin{equation}
S_{\mathrm{grad}}(W_\ell, q) = \operatorname{mean}\left|\frac{\partial q}{\partial W_\ell}\right|.
\end{equation}
The layer with maximal $S_{\mathrm{grad}}$ is assigned to $n$.

\textbf{Parameters (trend-based intervention).} For a physical parameter
with a known qualitative effect, we apply a controlled tensor intervention
$W_\ell' = (1+\epsilon)W_\ell$ (where the injected perturbation $\epsilon = \sigma|\mathcal{N}(0, 1)| \ge 0$ is strictly positive to preserve unambiguous directional trend evaluation), run a full forward pass, and measure the
relative change in the target observable,
\begin{equation}
S_{\mathrm{pert}} = \frac{|q(W_\ell') - q(W_\ell)|}{|q(W_\ell)| + \delta}.
\end{equation}
Where the parameter has a known physical direction, we additionally check
directional consistency (sign match with the analytically expected trend).
We use tensor intervention to discover neural-physical correspondences, while analytical derivatives of the physics residual with respect to physical parameters ($\partial \text{mean}(R^2)/\partial p$) provide an independent check of physical directional consistency.
For each node we compare the top score against the runner-up to assign a
confidence label (high/ambiguous/low), and report
$\text{Alignment Accuracy} = (\text{\# high-confidence nodes}/\text{total nodes})\times 100$,
with $\geq 90\%$ as our target threshold.

\subsection{Transfer Protocol}
To test whether discovered causal roles transfer beyond discovery conditions,
we define an in-distribution training domain $\mathcal{D}_{\text{ID}} = \{t \in [0, 1]\}$
and an out-of-distribution extrapolation domain $\mathcal{D}_{\text{OOD}} = \{t \in [1, 2]\}$.
The protocol is:
\begin{enumerate}
    \item Discover the optimal node-to-layer mapping $\mathcal{M}_{\text{ID}}: V \to \mathcal{L}$ using ID data.
    \item \emph{Freeze} the mapping $\mathcal{M}_{\text{ID}}$ without re-optimizing or re-searching on OOD data.
    \item Evaluate predictive error ($L_2$), physics residual adherence ($\text{mean}(R^2)$), and intervention response on $\mathcal{D}_{\text{OOD}}$ under the frozen mapping.
\end{enumerate}
This isolates whether internal representations maintain decoupled functional roles across distribution shifts, rather than merely fitting localized statistical correlations.

\section{Experimental Setup}
\label{sec:setup}
We evaluate GIIT on the 1D damped harmonic oscillator
($m\ddot{x}+c\dot{x}+kx=0$), modeled with a fully connected PINN
(Linear(1$\to$32)$\to\tanh\to$Linear(32$\to$32)$\to\tanh\to$Linear(32$\to$1), 2{,}209
parameters), trained with a combined data-fit and physics-residual loss.
We report three independent axes: predictive accuracy (relative $L_2$
error, target $\leq5\%$), physics adherence (mean residual norm), and
mechanistic consistency (sign/magnitude match of intervention response
against analytically expected sensitivity). The ID regime covers $t \in [0, 1]$ (500 collocation points),
while the OOD transfer regime tests temporal extrapolation over $t \in [1, 2]$.

\section{Main Results}
\label{sec:results}

\begin{table}[h]
\caption{Oscillator results across In-Distribution ($t\in[0,1]$) and Out-Of-Distribution ($t\in[1,2]$) regimes.}
\label{tab:main}
\centering
\small
\begin{tabular}{lrrl}
\toprule
Regime & Pred.\ $L_2$ error & Physics residual & Mech.\ consistency \\
\midrule
ID ($t \in [0, 1]$)  & $1.13\%$ & $2.007\times10^{2}$ &  (Passes $L_2 \le 5\%$, mapping identified) \\
OOD ($t \in [1, 2]$) & $2007.34\%$ & $1.317\times10^{5}$ &  (Severe decoupling, sign inversion) \\
\bottomrule
\end{tabular}
\end{table}

In-distribution, the model accurately predicts displacement ($L_2 = 1.13\%$) and maintains consistent sensitivity:
increasing mass slows decay ($\frac{\partial R}{\partial m} > 0$), increasing damping accelerates decay ($\frac{\partial R}{\partial c} < 0$).
In the representative run, perturbation mapping identifies hidden tensors associated with the physical parameters (\texttt{fch.1.0.weight} and \texttt{fch.0.0.weight}), while the multi-seed analysis (detailed in the appendix) reveals the underlying representation distribution.

\section{Transfer Analysis}
\label{sec:transfer}
Evaluating the frozen ID mapping under the OOD regime ($t \in [1, 2]$) reveals significant mechanistic divergence.
As shown in Table~\ref{tab:main}, while the model achieves low predictive error in-distribution, its OOD predictive error explodes ($L_2 = 2007.34\%$), and the physics residual norm increases by nearly three orders of magnitude ($1.317\times10^5$).
More critically, internal sensitivity signs invert under extrapolation: $\frac{\partial R}{\partial m}$ becomes negative ($-1747.92$) and $\frac{\partial R}{\partial c}$ becomes positive ($+579.62$), directly violating governing physical intuition. 
Furthermore, re-running auto-mapping in the OOD regime demonstrates that mass and damping mappings collapse to a shared early-layer tensor (\texttt{fch.0.0.weight}), while the spring-constant mapping shifts to a neighboring early layer (\texttt{fch.1.0.weight}).
This directly answers our central question: learned physical correspondences do \emph{not} automatically transfer OOD, demonstrating that PINNs can satisfy numerical loss objectives in-distribution while failing to form robust, generalizable internal representations.

\section{Discussion and Limitations}
\label{sec:limitations}
The oscillator experiments show GIIT successfully uncovering internal functional mappings in-distribution,
while our transfer protocol demonstrates how distributional shifts break internal mechanistic consistency before standard output metrics can diagnose the failure.
Limitations: the mapping is behavioral rather than a formal proof of causal identity; experiments use known governing equations and controlled, low-dimensional systems; and tensor representations may be entangled across complex multi-layer architectures.

\section{Conclusion}
\label{sec:conclusion}
GIIT provides an operational, model-agnostic complement to correlational
explainability tools for SciML: it maps physics-graph nodes to internal
tensors and tests the mapping under intervention. Our ID$\to$OOD transfer analysis reveals that
in-distribution constraint satisfaction does not ensure out-of-distribution mechanistic integrity,
establishing GIIT as a valuable diagnostic tool for verifying causal learning in scientific machine learning.

\bibliographystyle{plainnat}
\bibliography{references}

\newpage
\appendix
\section{Technical Appendix}
\label{sec:appendix}

\subsection{Seed Sweep and Representation Stability Analysis}
To rigorously evaluate whether the discovered neural-physical mappings are artifacts of a specific random noise draw or represent stable topological properties of the trained PINN, we executed a systematic seed sweep across 21 independent random perturbation seeds ($\text{seed} \in [0, 20]$ with noise std $\sigma = 0.1$).

Table~\ref{tab:seed_sweep} summarizes the modal (top) mapped neural tensor, the number of seeds where that tensor was selected as the highest-scoring component, and the distribution across candidate layers.

\begin{table}[h]
\caption{Neural Component Mapping Distributions across 21 Random Seeds ($\text{seed} \in [0, 20]$, $\sigma = 0.1$) for both In-Distribution ($t\in[0,1]$) and Out-Of-Distribution ($t\in[1,2]$) evaluation.}
\label{tab:seed_sweep}
\centering
\small
\begin{tabular}{lcccc}
\toprule
Physical Variable & \multicolumn{2}{c}{In-Distribution ($t \in [0, 1]$)} & \multicolumn{2}{c}{Out-Of-Distribution ($t \in [1, 2]$)} \\
\cmidrule(lr){2-3} \cmidrule(lr){4-5}
& Top Component & Consistency & Top Component & Consistency \\
\midrule
Mass ($m$)                  & \texttt{fch.1.0.weight} & 10/21 (47.6\%) & \texttt{fch.0.0.weight} & 16/21 (76.2\%) \\
Damping Coeff ($c$)         & \texttt{fch.1.0.weight} & 10/21 (47.6\%) & \texttt{fch.0.0.weight} & 16/21 (76.2\%) \\
Spring Constant ($k$)       & \texttt{fch.0.0.weight} & 8/21 (38.1\%)  & \texttt{fch.1.0.weight} & 12/21 (57.1\%) \\
Displacement ($x$)          & \texttt{fch.1.0.weight} & 11/21 (52.4\%) & \texttt{fch.1.0.weight} & 16/21 (76.2\%) \\
Velocity ($\dot{x}$)        & \texttt{fch.0.0.weight} & 17/21 (81.0\%) & \texttt{fch.0.0.weight} & 13/21 (61.9\%) \\
Acceleration ($\ddot{x}$)   & \texttt{fch.0.0.weight} & 16/21 (76.2\%) & \texttt{fch.0.0.weight} & 10/21 (47.6\%) \\
\bottomrule
\end{tabular}
\end{table}

\paragraph{Key Observations from Multi-Seed Sweep:}
\begin{itemize}
    \item \textbf{Derivative Localization Stability}: Physical derivative states exhibit high consistency in-distribution. Velocity ($\dot{x}$) and Acceleration ($\ddot{x}$) map to the first hidden block (\texttt{fch.0.0.weight}) in 81.0\% and 76.2\% of seeds, respectively.
    \item \textbf{OOD Representation Reorganization}: Under out-of-distribution temporal extrapolation ($t \in [1, 2]$), mass and damping mappings collapse to a shared early-layer tensor (\texttt{fch.0.0.weight} in 76.2\% of seeds), while the spring-constant mapping shifts to a neighboring early layer (\texttt{fch.1.0.weight} in 57.1\% of seeds), demonstrating that the network's internal capacity to distinguish independent physical parameters degrades outside the training envelope.
\end{itemize}

\subsection{Analytical Physics Residual Sensitivities}
Table~\ref{tab:sensitivities} provides the analytical gradients $\partial \text{mean}(R^2)/\partial p$ evaluated across both regimes, demonstrating the sign inversions discussed in Section~\ref{sec:transfer}.

\begin{table}[h]
\caption{Physics Residual Parameter Sensitivities ($\partial \text{mean}(R^2)/\partial p$).}
\label{tab:sensitivities}
\centering
\small
\begin{tabular}{lccc}
\toprule
Physical Parameter & Expected Sign & ID ($t \in [0, 1]$) & OOD ($t \in [1, 2]$) \\
\midrule
Mass ($m$)                  & $+$ & $+1377.28$ & $-1747.92$ (\textbf{Inverted}) \\
Damping Coeff ($c$)         & $-$ & $-3.84$    & $+579.62$ (\textbf{Inverted}) \\
Spring Constant ($k$)       & $+$ & $+111.77$  & $+458.12$ \\
\bottomrule
\end{tabular}
\end{table}

\subsection{Supplementary Material}
Anonymized code and ruleset specifications for the viscous Burgers' equation system are provided in a \texttt{.zip} archive as supplementary material.

\end{document}



